%% ============================================================
%% CHAPTER 5 — Characterization: What Bare Metal Costs and Buys
%%                                        (budget: ~8,000 words)
%% Role: the measurement chapter; answers the quantitative half of RQ1.
%% ============================================================
\chapter{Characterization: What Bare Metal Costs and Buys}
\label{ch:characterization}

\section{Static Footprint}
\label{sec:char-footprint}
% Measured flash/SRAM/PSRAM against nanobot-on-Pi (RSS, storage) and
% zclaw where public. Boot-to-ready time. Table + one bar figure.

\section{Per-Turn Latency Decomposition}
\label{sec:char-latency}
% The waterfall figure: DNS / TLS handshake / TTFB / streaming / parse /
% tool execution, vs the host reference loop on the same network (type-ii
% baseline isolates the MCU tax). TLS connection-reuse on/off. Findings
% paragraph: where the time actually goes.

\section{Energy}
\label{sec:char-energy}
% Joules per task (INA226 trace aligned via GPIO marker); idle draw;
% 24-hour always-on total vs Pi Zero 2W + nanobot. The battery-life /
% annual-energy projection figure — the headline chart of the chapter.

\section{Task Parity with Linux Runtimes}
\label{sec:char-parity}
% The two-table design: (1) parity table on the common task set (success
% with CIs, cost, latency, energy; success is the equivalence check, not
% the differentiator); (2) capability coverage matrix (host-only rows:
% shell, browser; MCU-only rows: GPIO, physical actuation). Model-tier
% sweep subsection: does harness sensitivity fall with stronger models;
% do deterministic guards compensate weaker ones.

\section{Reliability Under Faults}
\label{sec:char-faults}
% Fault-injection matrix results (WiFi outage 10/60/600s; mid-task power
% pull; API 429/500/malformed; SPIFFS full): recovery rate + MTTR, vs
% nanobot-on-Pi under the same faults AND vs MimiClaw with recovery
% mechanisms disabled (proves the design, not the SDK, earns the result).
% 7-day soak: heap-drift curve, crash/reconnect/loss counts.

\section{Architecture Validation}
\label{sec:char-dualcore}
% Dual-core split vs single-core pinning: channel-latency jitter during
% agent turns. Turns the §3.3 claim into a distribution plot.

\section{Discussion}
\label{sec:char-discussion}
% What parity + two-orders-of-magnitude resource gap means; where the MCU
% genuinely loses (flexibility, coverage); threats to validity.
